\section{Derivation of the nonlinear twist action}

\label{app:E}

For the fixed-scalar, round-shape family of Section 4.2 in the unrotated coframe the induced spatial metric, lapse and shift in the invariant basis are

\begin{equation}\tag{E1}\label{eq:E1}
\begin{aligned}
\gamma&=r^2I+L^2\omega\omega^{\mathsf T},&
\mathcal N&=\frac{Lr}{\sqrt{r^2+L^2\omega^2}},\\
\beta^i&=\frac{L^2\omega_i}{r^2+L^2\omega^2},&
\gamma^{-1}&=\frac1{r^2}\left(I-\frac{L^2\omega\omega^{\mathsf T}}{r^2+L^2\omega^2}\right).
\end{aligned}
\end{equation}

At each circle coordinate, the spatial shift is along the symmetry axis of $\gamma$ and its spatial Lie derivative of $\gamma$ vanishes. Thus $K=\gamma'/(2\mathcal N)$ up to an irrelevant overall sign convention for the normal. Direct matrix traces give

\begin{equation}\tag{E2}\label{eq:E2}
\begin{aligned}
\mathrm{tr}\!\left[(\gamma^{-1}K)^2\right]
-\left[\mathrm{tr}(\gamma^{-1}K)\right]^2
&=\frac{L^2}{2r^4}Q,\\
{}^{(3)}R&=\frac{24}{r^2}-\frac{8L^2}{r^4}\omega^2.
\end{aligned}
\end{equation}

The Euclidean curvature decomposition adds the negative of the first line and the divergence $-2\partial_y(\omega\cdot\omega')/r^2$ to $^{(3)}R$, giving (V2). Equations (A1)--(A2) give the same result without a metric decomposition. This comparison fixes both the sign and the normalization of the quartic term.

For an arbitrary smooth mixing rotation, let $E_3^y$ denote the circle component of the final frame vector dual to $e'^3$. With fixed $r,L,\eta,V$ the torsion decomposition gives

\begin{equation}\tag{E3}\label{eq:E3}
\begin{aligned}
R_{\mathrm{EC}}&=R_{\mathrm{LC}}-\frac{6\eta^2}{r^2}-\frac{2V^2}{3}
+2V\partial_yE_3^y,\\
N&=-\frac{6\eta}{r}\partial_yE_3^y.
\end{aligned}
\end{equation}

The volume density is constant and the spatial invariant divergence is zero. The last terms therefore integrate to zero for periodic profiles. For a single 03 rotation and $\omega=(w,0,0)$, $E_3^y=\cos\chi/L+w\sin\chi/r$, which exhibits the boundary current explicitly. This proves scalar-action independence from mixing in the circle family.

For completeness, with $p=\omega'$, $s=\omega\cdot p$, the derivative momentum and Euler expression are

\begin{equation}\tag{E4}\label{eq:E4}
\begin{aligned}
\frac{\partial\mathcal L}{\partial p}
&=\frac{C_\omega}{8}(\omega^2p-s\omega),\\
E&=2C_\omega\omega+\frac{C_\omega}{8}
\left[(2p^2+\omega\cdot\omega'')\omega
-2sp-\omega^2\omega''\right].
\end{aligned}
\end{equation}

Contracting with $\omega$ cancels the second derivatives and gives (V5). This is a pointwise proof, valid also where the polar coordinate $\rho,n$ is singular. It requires no division by the amplitude and establishes the zero-twist stationary result for all real smooth twist profiles in this fixed-scalar, round-shape family.
